\ifdefined\gkver\else\def\gkver{student}\fi
\documentclass[\gkver]{gaokaozhenti}
\begin{document}

\chapter{2009年广东理综}
%% paperType: 理综物理部分

\begin{enumerate}
\item
%% number: 1
%% paperName: 2009年广东理综第 1 题 4 分
%% typeId: 1
%% type: 单选题
%% chapter: 其他
%% point: 物理学史
%% method: 无
%% score: 4
%% degree: 650
%% duplicateId: 0
%% body:
物理学的发展丰富了人类对物质世界的认识，推动了科学技术的创新和革命，促进了物质生产的繁荣与人类文明的进步。下列表述正确的是\xzanswer{A}

\fourchoices[answer=A]
{牛顿发现了万有引力定律}
{光电效应证实了的光的波动性}
{洛伦兹发现了电磁感应定律}
{相对论的创立表明了经典力学已不再适用}

%% answer: A
%% memo:
\memoanswer{
电磁感应定律是法拉第发现的，B 错误；光电效应证实了光的粒子性，C 错误；相对论和经典力学研究的领域不同，不能说相对论的创立表明经典力学已不再适用，D 错误。正确答案选 A。
}

\item
%% number: 2
%% paperName: 2009年广东理综第 2 题 4 分
%% typeId: 1
%% type: 单选题
%% chapter: 原子和原子核
%% point: 核聚变
%% method: 无
%% score: 4
%% degree: 650
%% duplicateId: 0
%% body:
科学家发现在月球上含有丰富的 $\ce{^3_2 He}$（氦 3），它是一种高效、清洁、安全的核聚变燃料，其参与的一种核聚变反应的方程式为 $\ce{^3_2 He + ^3_2 He -> 2 ^1_1 H + ^4_2 He}$。关于 $\ce{^3_2 He}$ 聚变下列表述正确的是\xzanswer{B}

\fourchoices[answer=B]
{聚变反应不会释放能量}
{聚变反应产生了新的原子核}
{聚变反应没有质量亏损}
{目前核电站都采用 $\ce{^3_2 He}$ 聚变反应发电}

%% answer: B
%% memo:
\memoanswer{
聚变反应时将质量较小的轻核聚变成质量较大的核，聚变过程会有质量亏损，要放出大量的能量。但目前核电站都采用铀核的裂变反应。因此 B 正确。
}

\item
%% number: 3
%% paperName: 2009年广东理综第 3 题 4 分
%% typeId: 1
%% type: 单选题
%% chapter: 直线运动
%% point: 运动图象
%% method: 无
%% score: 4
%% degree: 650
%% duplicateId: 0
%% body:
某物体运动的速度图象如图所示，根据图象可知\xzanswer{AC}

\onepicture[w=5.5cm]{figs/03.png}

\fourchoices[answer=AC]
{$0\sim2\Us$ 内的加速度为 $1\Umsq$}
{$0\sim5\Us$ 内的位移为 $10\Um$}
{第 $1\Us$ 末与第 $3\Us$ 末的速度方向相同}
{第 $1\Us$ 末与第 $5\Us$ 末加速度方向相同}

%% answer: AC
%% memo:
\memoanswer{
$v-t$ 图象反映的是速度 $v$ 随时间 $t$ 的变化规律，其斜率表示的是加速度，A 正确；图中图象与坐标轴所围成的梯形面积表示的是 $0\sim5\Us$ 内的位移为 $7\Um$，B 错误。在前 $5\Us$ 内物体的速度都大于零，即运动方向相同，C 正确；$0\sim2\Us$ 加速度为正，$4\sim5\Us$ 加速度为负，方向不同，D 错误。正确选项：AC。
}

\item
%% number: 4
%% paperName: 2009年广东理综第 4 题 4 分
%% typeId: 1
%% type: 单选题
%% chapter: 光学
%% point: 光电效应
%% method: 无
%% score: 4
%% degree: 650
%% duplicateId: 0
%% body:
硅光电池是利用光电效应原理制成的器件。下列表述正确的是\xzanswer{A}

\fourchoices[answer=A]
{硅光电池是把光能转变为电能的一种装置}
{硅光电池中吸收了光子能量的电子都能逸出}
{逸出的光电子的最大初动能与入射光的频率无关}
{任意频率的光照射到硅光电池上都能产生光电效应}

%% answer: A
%% memo:
\memoanswer{
硅光电池是把光能转变为电能的一种装置，A 正确；是利用光电效应原理制成的器件，依据光电效应方程 $E_k=h\nu-W=h\nu-h\nu_0$ 可见只有当入射光子的频率大于极限频率时才可能发生光电效应，B 错误，C 错误，D 错误。
}

\item
%% number: 5
%% paperName: 2009年广东理综第 5 题 2 分
%% typeId: 1
%% type: 单选题
%% chapter: 曲线运动
%% point: 人造地球卫星
%% method: 无
%% score: 2
%% degree: 650
%% duplicateId: 0
%% body:
发射人造卫星是将卫星以一定的速度送入预定轨道。发射场一般选择在尽可能靠近赤道的地方。如图所示，这样选址的优点是，在赤道附近\xzanswer{B}

\onepicture[w=3.2cm]{figs/05.png}

\fourchoices[answer=B]
{地球的引力较大}
{地球自转线速度较大}
{重力加速度较大}
{地球自转角速度较大}

%% answer: B
%% memo:
\memoanswer{
由于发射卫星需要将卫星以一定的速度送入运动轨道，在靠近赤道处的地面上的物体的线速度最大，发射时较节能，因此 B 正确。
}

\item
%% number: 6
%% paperName: 2009年广东理综第 6 题 4 分
%% typeId: 1
%% type: 单选题
%% chapter: 电场
%% point: 带电粒子的直线运动
%% method: 无
%% score: 4
%% degree: 450
%% duplicateId: 0
%% body:
如图所示，在一个粗糙水平面上，彼此靠近地放置两个带同种电荷的小物块。由静止释放后，两个物块向相反方向运动，并最终停止。在物块的运动过程中，下列表述正确的是\xzanswer{A}

\onepicture[w=4.8cm]{figs/06.png}

\fourchoices[answer=A]
{两个物块的电势能逐渐减少}
{物块受到的库仑力不做功}
{两个物块的机械能守恒}
{物块受到的摩擦力始终小于其受到的库仑力}

%% answer: A
%% memo:
\memoanswer{
由于两电荷电性相同，则二者之间的作用力为斥力，因此在远离过程中，电场力做正功，则电势能逐渐减少，A 正确；B 错误；由于运动过程中，有重力以外的力电场力和摩擦力做功，故机械能不守恒，C 错误；在远离过程中开始电场力大于摩擦力，后来电场力小于摩擦力。
}

\item
%% number: 7
%% paperName: 2009年广东理综第 7 题 4 分
%% typeId: 2
%% type: 多选题
%% chapter: 力物体的平衡
%% point: 弹力
%% method: 无
%% score: 4
%% degree: 450
%% duplicateId: 0
%% body:
某缓冲装置可抽象成如图所示的简单模型。图中 $k_1$、$k_2$ 为原长相等、劲度系数不同的轻质弹簧。下列表述正确的是\xzanswer{BD}

\onepicture[w=4.8cm]{figs/07.png}

\fourchoices[answer=BD]
{缓冲效果与弹簧的劲度系数无关}
{垫片向右移动时，两弹簧产生的弹力大小相等}
{垫片向右移动时，两弹簧的长度保持相等}
{垫片向右移动时，两弹簧的弹性势能发生改变}

%% answer: BD
%% memo:
\memoanswer{
不同弹簧的缓冲效果与弹簧的劲度系数有关，A 错误；在垫片向右运动的过程中，由于两个弹簧相连，则它们之间的作用力等大，B 正确；由于两弹簧的劲度系数不同，由胡克定律 $F=k\Delta x$ 可知，两弹簧的形变量不同，则两弹簧的长度不相等，C 错误；在垫片向右运动的过程中，由于弹簧的弹力做功，则弹性势能将发生变化，D 正确。
}

\item
%% number: 8
%% paperName: 2009年广东理综第 8 题 4 分
%% typeId: 2
%% type: 多选题
%% chapter: 运动和力
%% point: 超重失重
%% method: 无
%% score: 4
%% degree: 650
%% duplicateId: 0
%% body:
某人在地面上用弹簧秤称得其体重为 $490\UN$，他将弹簧秤移至电梯内称其体重，$t_0$ 至 $t_1$ 时间段内，弹簧秤的示数如图所示，电梯运行的 $v-t$ 图象可能是（取电梯向上运动的方向为正）\xzanswer{AD}

\onepicture[w=5.0cm]{figs/08a.png}

\fourchoices[answer=AD, ispicture=true, h=1.9cm]
{figs/08b.png}
{figs/08c.png}
{figs/08d.png}
{figs/08e.png}

%% answer: AD
%% memo:
\memoanswer{
由图可知，在 $t_0\sim t_1$ 时间内，弹簧秤的示数小于实际重量，则处于失重状态，此时具有向下的加速度，在 $t_1\sim t_2$ 阶段弹簧秤示数等于实际重量，则既不超重也不失重，在 $t_2\sim t_3$ 阶段，弹簧秤示数大于实际重量，则处于超重状态，具有向上的加速度，若电梯向下运动，则 $t_0\sim t_1$ 时间内向下加速，$t_1\sim t_2$ 阶段匀速运动，$t_2\sim t_3$ 阶段减速下降，A 正确；同理 D 项也符合题意；B、C 项 $t_0\sim t_1$ 内超重，不符合题意。
}

\item
%% number: 9
%% paperName: 2009年广东理综第 9 题 4 分
%% typeId: 2
%% type: 多选题
%% chapter: 电磁感应
%% point: 远距离输电
%% method: 无
%% score: 4
%% degree: 450
%% duplicateId: 0
%% body:
如图所示为远距离高压输电的示意图。关于远距离输电，下列表述正确的是\xzanswer{ABD}

\onepicture[w=6.0cm]{figs/09.png}

\fourchoices[answer=ABD]
{增加输电导线的横截面积有利于减少输电过程中的电能损失}
{高压输电是通过减小输电电流来减小电路的发热损耗}
{在输送电压一定时，输送的电功率越大，输电过程中的电能损失越小}
{高压输电必须综合考虑各种因素，不一定是电压越高越好}

%% answer: ABD
%% memo:
\memoanswer{
依据输电原理，电路中的功率损耗 $\Delta P=I^2R_{\text{线}}$，而 $R_{\text{线}}=\rho\frac{L}{S}$，增大输电线的横截面积，减小输电线的电阻，则能够减小输电线上的功率损耗，A 正确；

由 $P=UI$ 来看在输送功率一定的情况下，输送电压 $U$ 越大，则输电电流越小，则功率损耗越小，B 正确；

若输电电压一定，输送功率越大，则电流 $I$ 越大，电路中损耗的电功率越大，C 错误；

输电电压并不是电压越高越好，因为电压越高，对于安全和技术的要求越高，因此并不是输电电压越高越好，D 正确。
}

\item
%% number: 10
%% paperName: 2009年广东理综第 10 题 4 分
%% typeId: 1
%% type: 单选题
%% chapter: 电路
%% point: 动态分析
%% method: 无
%% score: 4
%% degree: 450
%% duplicateId: 0
%% body:
如图所示，电动势为 $E$、内阻不计的电源与三个灯泡和三个电阻相接。只合上开关 $S_1$，三个灯泡都能正常工作。如果再合上 $S_2$，则下列表述正确的是\xzanswer{C}

\onepicture[w=5.5cm]{figs/10.png}

\fourchoices[answer=C]
{电源输出功率减小}
{$L_1$ 上消耗的功率增大}
{通过 $R_1$ 上的电流增大}
{通过 $R_3$ 上的电流增大}

%% answer: C
%% memo:
\memoanswer{
在合上 $S_2$ 之前，三灯泡都能正常工作，合上 $S_2$ 之后，电路中的总电阻 $R_{\text{总}}$ 减小，则 $I_{\text{总}}$ 增大，即流过 $R_1$ 的电流增大，由于不计内阻，电源的输出功率 $P_{\text{出}}=EI$，可见电源的输出功率增大，A 错误；$R_1$ 两端的电压增大，则并联部分的电压减小，$I_4$ 减小，$I_2$ 减小，$I_1$ 减小，可见 C 正确。
}

\item
%% number: 11
%% paperName: 2009年广东理综第 11 题 4 分
%% typeId: 2
%% type: 多选题
%% chapter: 磁场
%% point: 带电粒子在磁场中的运动
%% method: 无
%% score: 4
%% degree: 450
%% duplicateId: 0
%% body:
如图所示，表面粗糙的斜面固定于地面上，并处于方向垂直纸面向外、强度为 $B$ 的匀强磁场中。质量为 $m$、带电量为 $+Q$ 的小滑块从斜面顶端由静止下滑。在滑块下滑的过程中，下列判断正确的是\xzanswer{CD}

\onepicture[w=4.8cm]{figs/11.png}

\fourchoices[answer=CD]
{滑块受到的摩擦力不变}
{滑块到达地面时的动能与 $B$ 的大小无关}
{滑块受到的洛伦兹力方向垂直斜面向下}
{$B$ 很大时，滑块可能静止于斜面上}

%% answer: CD
%% memo:
\memoanswer{
取物块为研究对象，小滑块沿斜面下滑由于受到洛伦兹力作用，如图所示，C 正确；$N=mg\cos\theta+qvB$，由于 $v$ 不断增大，则 $N$ 不断增大，滑动摩擦力 $f=\mu N$，摩擦力增大，A 错误；滑块的摩擦力与 $B$ 有关，摩擦力做功与 $B$ 有关，依据动能定理，在滑块下滑到地面的过程中，满足 $\frac{1}{2}mv^2-0=mgh-fs$，所以滑块到地面时的动能与 $B$ 有关，B 错误；当 $B$ 很大，则摩擦力有可能很大，所以滑块可能静止在斜面上，D 正确。

\onepicture[w=5.0cm]{figs/11_an.jpg}
}

\item
%% number: 12
%% paperName: 2009年广东理综第 12 题 4 分
%% typeId: 2
%% type: 多选题
%% chapter: 磁场
%% point: 带电粒子在复合场中的运动
%% method: 无
%% score: 4
%% degree: 450
%% duplicateId: 0
%% body:
如图所示是质谱仪的工作原理示意图。带电粒子被加速电场加速后，进入速度选择器。速度选择器内相互正交的匀强磁场和匀强电场的强度分别为 $B$ 和 $E$。平板 S 上有可让粒子通过的狭缝 P 和记录粒子位置的胶片 $A_1A_2$。平板 S 下方有强度为 $B_0$ 的匀强磁场。下列表述正确的是\xzanswer{ABC}

\onepicture[w=4.2cm]{figs/12.png}

\fourchoices[answer=ABC]
{质谱仪是分析同位素的重要工具}
{速度选择器中的磁场方向垂直纸面向外}
{能通过狭缝 P 的带电粒子的速率等于 $E/B$}
{粒子打在胶片上的位置越靠近狭缝 P，粒子的荷质比越小}

%% answer: ABC
%% memo:
\memoanswer{
由加速电场可见粒子所受电场力向下，即粒子带正电，在速度选择器中，电场力水平向右，洛伦兹力水平向左，如图所示，因此速度选择器中磁场方向垂直纸面向外 B 正确；经过速度选择器时满足 $qE=qvB$，可知能通过的狭缝 P 的带电粒子的速率等于 $E/B$，带电粒子进入磁场做匀速圆周运动则有 $R=\frac{mv}{qB}$，可见当 $v$ 相同时，$R\propto\frac{m}{q}$，所以可以用来区分同位素，且 $R$ 越大，比荷就越大，D 错误。
}

\item
%% number: 13
%% paperName: 2009年广东理综第 13 题 6 分
%% typeId: 3
%% type: 填空题
%% chapter: 气体性质
%% point: 理想气体状态方程
%% method: 无
%% score: 6
%% degree: 650
%% duplicateId: 0
%% body:
\begin{enumerate}
	\item
	远古时代，取火是一件困难的事，火一般产生于雷击或磷的自燃。随着人类文明的进步，出现了“钻木取火”等方法。“钻木取火”是通过 \tkanswer[2cm]{} 方式改变物体的内能，把 \tkanswer[2cm]{} 转变成内能。
	\item
	某同学做了一个小实验：先把空的烧瓶放入冰箱冷冻，一小时后取出烧瓶，并迅速把一个气球紧密地套在瓶颈上，然后将烧瓶放进盛满热水的烧杯里，气球逐渐膨胀起来，如图所示。这是因为烧瓶里的气体吸收了水的 \tkanswer[2cm]{}，温度 \tkanswer[2cm]{}，体积 \tkanswer[2cm]{}。
\end{enumerate}
\onepicture[w=4.2cm, align=right]{figs/13.png}

%% answer:
\jdanswer{
\begin{enumerate}
	\item
	做功；机械能
	\item
	热量；升高；增大
\end{enumerate}
}

%% memo:
\memoanswer{
当用气球封住烧瓶，在瓶内就封闭了一定质量的气体，当将瓶子放到热水中，瓶内气体将吸收水的热量，增加气体的内能，温度升高，由理想气体状态方程 $\frac{pV}{T}=C$ 可知，气体体积增大。
}

\item
%% number: 14
%% paperName: 2009年广东理综第 14 题 6 分
%% typeId: 3
%% type: 填空题
%% chapter: 机械波
%% point: 干涉衍射
%% method: 无
%% score: 6
%% degree: 650
%% duplicateId: 0
%% body:
\begin{enumerate}
	\item
	在阳光照射下，充满雾气的瀑布上方常常会出现美丽的彩虹。彩虹是太阳光射入球形水珠经折射、内反射，再折射后形成的。光的折射发生在两种不同介质的 \tkanswer[2cm]{} 上，不同的单色光在同种均匀介质中 \tkanswer[2cm]{} 不同。
	\item
	如图所示为声波干涉演示仪的原理图。两个 U 形管 A 和 B 套在一起，A 管两侧各有一小孔。声波从左侧小孔传入管内，被分成两列频率 \tkanswer[2cm]{} 的波。当声波分别通过 A、B 传播到右侧小孔时，若两列波传播的路程相差半个波长，则此处声波的振幅 \tkanswer[3cm]{}；若传播的路程相差一个波长，则此处声波的振幅 \tkanswer[3cm]{}。
\end{enumerate}
\onepicture[w=4.2cm, align=right]{figs/14.png}

%% answer:
\jdanswer{
\begin{enumerate}
	\item
	界面；传播速度
	\item
	相等；等于零；等于原振幅的 2 倍
\end{enumerate}
}

%% memo:
\memoanswer{
声波从左侧小孔传入管内向上向下分别形成两列频率相同的波，若两列波传播的路程相差半个波长，则振动相消，所以此处振幅为零，若传播的路程相差一个波长，振动加强，则此处声波的振幅为原振幅的二倍。
}

\item
%% number: 15
%% paperName: 2009年广东理综第 15 题 10 分
%% typeId: 4
%% type: 实验题
%% chapter: 功和能
%% point: 验证动能定理
%% method: 无
%% score: 10
%% degree: 450
%% duplicateId: 0
%% body:
某实验小组利用拉力传感器和速度传感器探究“动能定理”。如图所示，他们将拉力传感器固定在小车上，用不可伸长的细线将其通过一个定滑轮与钩码相连，用拉力传感器记录通过 A、B 时的速度大小。小车中可以放置砝码。

\begin{enumerate}
	\item
	实验主要步骤如下：
\begin{steps}
	\item
	测量 \tkanswer[2cm]{} 和拉力传感器的总质量 $M_1$；把细线的一端固定在拉力传感器上，另一端通过定滑轮与钩码相连；正确连接所需电路；
	\item
	将小车停在 C 点，\tkanswer[2cm]{}，小车在细线拉动下运动，记录细线拉力及小车通过 A、B 时的速度。
	\item
	在小车中增加砝码，或 \tkanswer[2cm]{}，重复②的操作。
\end{steps}
	\item
	表 1 是他们测得的一组数据，其中 $M$ 是 $M_1$ 与小车中砝码质量之和，$|v_2^2-v_1^2|$ 是两个速度传感器记录速度的平方差，可以据此计算出动能变化量 $\Delta E$，$F$ 是拉力传感器受到的拉力，$W$ 是 $F$ 在 A、B 间所作的功。表格中的 $\Delta E_3=$ \tkanswer[1.5cm]{}，$W_3=$ \tkanswer[1.5cm]{}。（结果保留三位有效数字）

\begin{center}
\begin{tabular}{|c|c|c|c|c|c|}
\hline
次数 & $M/\Ukg$ & $\left|v_2^2-v_1^2\right|(\Ums)^2$ & $\Delta E/\UJ$ & $F/\UN$ & $W/\UJ$\\\hline
1 & 0.500 & 0.760 & 0.190 & 0.400 & 0.200\\\hline
2 & 0.500 & 1.65 & 0.413 & 0.840 & 0.420\\\hline
3 & 0.500 & 2.40 & $\Delta E_3$ & 1.220 & $W_3$\\\hline
4 & 1.000 & 2.40 & 1.20 & 2.420 & 1.21\\\hline
5 & 1.000 & 2.84 & 1.42 & 2.860 & 1.43\\\hline
\end{tabular}
\end{center}

	\item
	根据表 1，请在图中的方格纸上作出 $\Delta E-W$ 图线。
\end{enumerate}
\twopicture[wA=6.5cm, wB=3.6cm, align=right]
{figs/15a.png}
{figs/15b.png}

%% answer:
\jdanswer{
\begin{enumerate}
	\item
	小车；接通电源后释放小车；减少钩码的个数
	\item
	$0.619$；$0.610$
	\item
	如图
\end{enumerate}
}

%% memo:
\memoanswer{
（1）略；（2）由各组数据可见规律 $\Delta E=\frac{1}{2}m\left|v_2^2-v_1^2\right|$，可得 $\triangle E_3=0.600$；观察 $F-W$ 数据规律可得数值上 $W=F/2=0.610$；（3）在方格纸上作出 $\Delta E-W$ 图线如图所示。

\onepicture[w=5.5cm]{figs/15_an.jpg}
}

\item
%% number: 16
%% paperName: 2009年广东理综第 16 题 14 分
%% typeId: 4
%% type: 实验题
%% chapter: 电路
%% point: 伏安法测电阻
%% method: 无
%% score: 14
%% degree: 450
%% duplicateId: 0
%% body:
某实验小组利用实验室提供的器材探究一种金属丝的电阻率。所用的器材包括：输出为 $3\UV$ 的直流稳压电源、电流表、待测金属丝、螺旋测微器（千分尺）、米尺、电阻箱、开关和导线等。

\begin{enumerate}
	\item
	他们截取了一段金属丝，拉直后固定在绝缘的米尺上，并在金属丝上夹上一个小金属夹，金属夹可在金属丝上移动。请根据现有器材，设计实验电路，并连接电路实物图 1。
	\onepicture[w=5.5cm, align=right]{figs/16a.png}
	\item
	实验的主要步骤如下：
\begin{steps}
	\item
	正确连接电路，设定电阻箱的阻值，开启电源，合上开关；
	\item
	读出电流表的示数，记录金属夹的位置；
	\item
	断开开关，\tkanswer[2.5cm]{}，合上开关，重复②的操作。
\end{steps}
	\item
	该小组测得电流与金属丝接入长度关系的数据，并据此绘出如图 2 所示的关系图线，其斜率为 \tkanswer[2cm]{} $\UA^{-1}\cdot\Umn$（保留三位有效数字）；图线纵轴截距与电源电动势的乘积代表了 \tkanswer[3cm]{} 的电阻之和。
	\onepicture[w=5.0cm, align=right]{figs/16b.png}
	\item
	他们使用螺旋测微器测量金属丝的直径，示数如图 3 所示。金属丝的直径是 \tkanswer[2cm]{}。图中图线的斜率、电源电动势和金属丝横截面积的乘积代表的物理量是 \tkanswer[2.5cm]{}，其数值和单位为 \tkanswer[3cm]{}（保留三位有效数字）。
	\onepicture[w=3.0cm, align=right]{figs/16c.png}
\end{enumerate}

%% answer:
\jdanswer{
\begin{enumerate}
	\item
	电路图与实物连线如图所示
	\item
	测出接入电路的金属丝的长度
	\item
	$10.7$；电源内电阻与电阻箱
	\item
	$0.200\Umm$；金属丝的电阻率；$1.54\times10^{-7}\UOm$
\end{enumerate}
}

%% memo:
\memoanswer{
依据实验器材和实验目的测量金属丝的电阻率，电路图如图所示；电路实物图如图所示，依据闭合电路欧姆定律得 $E=I(r+R_0+R_x)$，参照题目给出的图象可得 $\frac{1}{I}=\frac{r+R_0}{E}+\frac{\rho}{ES}\cdot L$，可见直线的斜率 $k=\frac{\rho}{ES}$，可知斜率、电源电动势和金属丝横截面积的乘积代表的物理量是金属的电阻率 $\rho$，其数值和单位为 $1.54\times10^{-7}\UOm$；依据直线可得其斜率为 $1.63\UA^{-1}\cdot\Umn$，截距为 $\frac{r+R_0}{E}$，则图线纵轴截距与电源电动势的乘积为 $(r+R_0)$；金属丝的直径是 $0.200\Umm$。

\onepicture[w=8.0cm]{figs/16_an.jpg}
}

\item
%% number: 17
%% paperName: 2009年广东理综第 17 题 6 分
%% typeId: 6
%% type: 计算题
%% chapter: 曲线运动
%% point: 圆周运动的应用
%% method: 无
%% score: 6
%% degree: 450
%% duplicateId: 0
%% body:
\begin{enumerate}
	\item
	为了清理堵塞河道的冰凌，空军实施投弹爆破。飞机在河道上空高 $H$ 处以速度 $v_0$ 水平匀速飞行，投掷下炸弹并击中目标。求炸弹刚脱离飞机到击中目标所飞行的水平距离及击中目标时的速度大小。（不计空气阻力）
	\item
	如图所示，一个竖直放置的圆锥筒可绕其中心轴 $OO'$ 转动，筒内壁粗糙，筒口半径和筒高分别为 $R$ 和 $H$，筒内壁 A 点的高度为筒高的一半。内壁上有一质量为 $m$ 的小物块。求
\begin{enumerate}
	\item
	当筒不转动时，物块静止在筒壁 A 点受到的摩擦力和支持力的大小；
	\item
	当物块在 A 点随筒做匀速转动，且其所受到的摩擦力为零时，筒转动的角速度。
\end{enumerate}
\end{enumerate}
\onepicture[w=5.0cm, align=right]{figs/17.png}

%% answer:
\jdanswer{
\begin{enumerate}
	\item
	$s=v_0\sqrt{\frac{2H}{g}}$；$v=\sqrt{2gH+v_0^2}$
	\item
	① $f=\frac{mgH}{\sqrt{R^2+H^2}}$，$N=\frac{mgR}{\sqrt{R^2+H^2}}$；② $\omega=\frac{\sqrt{2gH}}{R}$
\end{enumerate}
}

%% memo:
\memoanswer{
（1）设飞行的水平距离为 $s$，在竖直方向上 $H=\frac{1}{2}gt^2$，得飞行时间为 $t=\sqrt{\frac{2H}{g}}$，则飞行的水平距离为 $s=v_0t=v_0\sqrt{\frac{2H}{g}}$。设击中目标时的速度为 $v$，飞行过程中，由机械能守恒得 $mgH+\frac{1}{2}mv_0^2=\frac{1}{2}mv^2$，得击中目标时的速度为 $v=\sqrt{2gH+v_0^2}$。

（2）物块受力如图所示。①由平衡条件得 $N-mg\cos\theta=0$，$f-mg\sin\theta=0$，其中 $\sin\theta=\frac{H}{\sqrt{R^2+H^2}}$，得摩擦力为 $f=mg\sin\theta=\frac{mgH}{\sqrt{R^2+H^2}}$，支持力为 $N=mg\cos\theta=\frac{mgR}{\sqrt{R^2+H^2}}$。

\onepicture[w=4.5cm]{figs/17_an1.jpg}

②这时物块的受力如图所示，由牛顿第二定律得 $mg\tan\theta=ma=m\frac{R}{2}\omega^2$，得筒转动的角速度为 $\omega=\sqrt{\frac{2g\tan\theta}{R}}=\frac{\sqrt{2gH}}{R}$。

\onepicture[w=4.0cm]{figs/17_an2.jpg}
}

\item
%% number: 18
%% paperName: 2009年广东理综第 18 题 15 分
%% typeId: 6
%% type: 计算题
%% chapter: 电磁感应
%% point: 电磁感应电路分析
%% method: 无
%% score: 15
%% degree: 450
%% duplicateId: 0
%% body:
如图 \ref{2009广东理综18a} 所示，一个电阻值为 $R$，匝数为 $n$ 的圆形金属线圈与阻值为 $2R$ 的电阻 $R_1$ 连接成闭合回路。线圈的半径为 $r_1$。在线圈中半径为 $r_2$ 的圆形区域内存在垂直于线圈平面向里的匀强磁场，磁感应强度 $B$ 随时间 $t$ 变化的关系图线如图 \ref{2009广东理综18b} 所示。图线与横、纵轴的截距分别为 $t_0$ 和 $B_0$。导线的电阻不计。求 $0$ 至 $t_1$ 时间内

\begin{enumerate}
	\item
	通过电阻 $R_1$ 上的电流大小和方向；
	\item
	通过电阻 $R_1$ 上的电量 $q$ 及电阻 $R_1$ 上产生的热量。
\end{enumerate}
\twopicture[labelA=2009广东理综18a, labelB=2009广东理综18b, numA=（a）, numB=（b）, wA=4.5cm, wB=4.5cm, align=right]
{figs/18a.png}
{figs/18b.png}

%% answer:
\jdanswer{
\begin{enumerate}
	\item
	$I=\frac{n\pi B_0r_2^2}{3Rt_0}$，电流由 b 向 a 通过 $R_1$
	\item
	$q=\frac{\pi nB_0r_2^2t_1}{3Rt_0}$；$Q=\frac{2n^2\pi^2B_0^2r_2^4t_1}{9Rt_0^2}$
\end{enumerate}
}

%% memo:
\memoanswer{
（1）由法拉第电磁感应定律得感应电动势为 $E=n\frac{\Delta\Phi}{\Delta t}=n\pi r_2^2\frac{\Delta B}{\Delta t}=\frac{n\pi B_0r_2^2}{t_0}$，由闭合电路的欧姆定律，得通过 $R_1$ 的电流大小为 $I=\frac{E}{3R}=\frac{n\pi B_0r_2^2}{3Rt_0}$，由楞次定律知该电流由 b 向 a 通过 $R_1$。

（2）由 $I=\frac{q}{t}$ 得在 $0$ 至 $t_1$ 时间内通过 $R_1$ 的电量为 $q=It_1=\frac{n\pi B_0r_2^2t_1}{3Rt_0}$，由焦耳定律得在 $0$ 至 $t_1$ 时间内 $R_1$ 产生的热量为 $Q=I^2R_1t_1=\frac{2n^2\pi^2B_0^2r_2^4t_1}{9Rt_0^2}$。
}

\item
%% number: 19
%% paperName: 2009年广东理综第 19 题 16 分
%% typeId: 6
%% type: 计算题
%% chapter: 运动和力
%% point: 动量和能量
%% method: 分情况讨论
%% score: 16
%% degree: 450
%% duplicateId: 0
%% body:
如图所示，水平地面上静止放置着物块 B 和 C，相距 $l=1.0\Um$。物块 A 以速度 $v_0=10\Ums$ 沿水平方向与 B 正碰。碰撞后 A 和 B 牢固地粘在一起向右运动，并再与 C 发生正碰，碰后瞬间 C 的速度 $v=2.0\Ums$。已知 A 和 B 的质量均为 $m$，C 的质量为 A 质量的 $k$ 倍，物块与地面的动摩擦因数 $\mu=0.45$。（设碰撞时间很短，$g$ 取 $10\Umsq$）

\begin{enumerate}
	\item
	计算与 C 碰撞前瞬间 AB 的速度；
	\item
	根据 AB 与 C 的碰撞过程分析 $k$ 的取值范围，并讨论与 C 碰撞后 AB 的可能运动方向。
\end{enumerate}
\onepicture[w=6.0cm, align=right]{figs/19.png}

%% answer:
\jdanswer{
\begin{enumerate}
	\item
	$4\Ums$
	\item
	$2\leq k\leq6$；当 $2\leq k<4$ 时，$v_3>0$，即与 C 碰撞后，AB 向右运动；当 $k=4$ 时，$v_3=0$，即 AB 停止；当 $4<k\leq7.74$ 时，$v_3<0$，即 AB 向左运动
\end{enumerate}
}

%% memo:
\memoanswer{
（1）设 A、B 碰后速度为 $v_1$，由于碰撞时间很短，A、B 相碰的过程动量守恒，得 $mv_0=2mv_1$ ①。在 A、B 向 C 运动，设与 C 碰撞前速度为 $v_2$，在此过程中由动能定理，有 $-\mu 2mgl=\frac{1}{2}\cdot2mv_2^2-\frac{1}{2}\cdot2mv_1^2$ ②，得 A、B 与 C 碰撞前的速度为 $v_2=\sqrt{v_1^2-2\mu gl}=4\Ums$。

（2）设 A、B 与 C 碰后速度为 $v_3$，A、B 与 C 碰撞的过程动量守恒 $2mv_2=2mv_3+kmv$ ③，$v_3=\frac{2v_2-kv}{2}=(4-k)\Ums$ ④。碰后 A、B 的速度 $v_3$ 必须满足 $v_3\leq v$ ⑤，$\frac{1}{2}\cdot2mv_2^2\leq\frac{1}{2}\cdot2mv_3^2+\frac{1}{2}kmv^2$ ⑥。由④⑤⑥式得 $2\leq k\leq6$。由④式知：当 $2\leq k<4$ 时，$v_3>0$，即与 C 碰撞后，AB 向右运动；当 $k=4$ 时，$v_3=0$，即与 C 碰撞后，AB 停止；当 $4<k\leq7.74$ 时，$v_3<0$，即与 C 碰撞后，AB 向左运动。
}

\item
%% number: 20
%% paperName: 2009年广东理综第 20 题 17 分
%% typeId: 6
%% type: 计算题
%% chapter: 电场
%% point: 带电粒子的直线运动
%% method: 无
%% score: 17
%% degree: 250
%% duplicateId: 0
%% body:
如图所示，绝缘长方体 B 置于水平面上，两端固定一对平行带电极板，极板间形成匀强电场 $E$。长方体 B 的上表面光滑，下表面与水平面的动摩擦因数 $\mu=0.05$（设最大静摩擦力与滑动摩擦力相同）。B 与极板的总质量 $m_B=1.0\Ukg$。带正电的小滑块 A 质量 $m_A=0.60\Ukg$，其受到的电场力大小 $F=1.2\UN$。假设 A 所带的电量不影响极板间的电场分布。$t=0$ 时刻，小滑块 A 从 B 表面上的 a 点以相对地面的速度 $v_A=1.5\Ums$ 向左运动，同时，B（连同极板）以相对地面的速度 $v_B=0.40\Ums$ 向右运动。问（$g$ 取 $10\Umsq$）

\begin{enumerate}
	\item
	A 和 B 刚开始运动时的加速度大小分别为多少？
	\item
	若 A 最远能到达 b 点，a、b 的距离 $L$ 应为多少？从 $t=0$ 时刻至 A 运动到 b 点时，摩擦力对 B 做的功为多少？
\end{enumerate}
\onepicture[w=5.5cm, align=right]{figs/20.png}

%% answer:
\jdanswer{
\begin{enumerate}
	\item
	$2\Umsq$，方向水平向右；$2\Umsq$，方向水平向左
	\item
	$0.62\Um$，$-0.072\UJ$
\end{enumerate}
}

%% memo:
\memoanswer{
（1）A、B 刚开始运动时，由牛顿第二定律，得：对于 A，只受电场力，有 $F=m_Aa_A$，得 A 的加速度大小为 $a_A=\frac{F}{m_A}=2\Umsq$，方向水平向右；对于 B 及极板，由牛顿第三定律知所受电场力为 $F'=F=1.2\UN$，由牛顿第二定律得 $F'+f=m_Ba_B$，而 $f=\mu N=\mu(m_A+m_B)g$，得 B 的加速度大小为 $a_B=\frac{F+\mu(m_A+m_B)g}{m_B}=2\Umsq$，方向水平向左。

（2）A、B 开始时都做加速度为 $2\Umsq$ 的匀减速直线运动，而 $v_A>v_B$，故 B 先减速到 0，设从 $t=0$ 到 B 速度减为 0 的时间为 $t_1$，则 $v_B=a_Bt_1$，得 $t_1=\frac{v_B}{a_B}=0.2\Us$。在 $t_1$ 时间内，A 的向左运动的位移为 $s_A=v_At_1-\frac{1}{2}a_At_1^2=0.28\Um$，B 向右运动的位移为 $s_B=v_Bt_1-\frac{1}{2}a_Bt_1^2=0.04\Um$。$t_1$ 时，A 的速度为 $v_A'=v_A-a_At_1=1.2\Ums$。此后 A 继续向左做加速度为 $2\Umsq$ 的匀减速直线运动，而 B 受到向左的电场力 $F>f$，故 B 向左做匀加速运动，设经时间 $t_2$，A、B 速度达到相同的 $v$，则 $F'-f=m_Ba_B'$，得 $a_B'=\frac{F'-\mu(m_A+m_B)g}{m_B}=0.4\Umsq$，$v_A'-a_At_2=a_B't_2$，得 $t_2=0.5\Us$。在时间 $t_2$ 内，A 向左的位移为 $s_A'=v_A't_2-\frac{1}{2}a_At_2^2=0.35\Um$，B 向左的位移为 $s_B'=\frac{1}{2}a_B't_2^2=0.05\Um$，则 a、b 间距离为 $L=s_A+s_B+s_A'-s_B'=0.62\Um$。在此过程中摩擦力对 B 做的功为 $W=-\mu(m_A+m_B)g(s_B+s_B')=-0.072\UJ$。
}

\end{enumerate}

\end{document}
